\chapter{2015年天津卷（理）}

\section{单选题}

\begin{problem}
已知全集 \(U = \{1, 2, 3, 4, 5, 6, 7, 8\}\)，集合 \(A = \{2, 3, 5, 6\}\)，集合 \(B = \{1, 3, 4, 6, 7\}\)，则集合 \(A \cap \complement_U B =\)（\quad）

\choices
    {\(\{2,5\}\)}
    {\(\{3,6\}\)}
    {\(\{2,5,6\}\)}
    {\(\{2,3,5,6,8\}\)}
\end{problem}

\begin{problem}
设变量 \(x\)，\(y\) 满足约束条件 \(\begin{cases}
x + 2 \geqslant 0,\\
x - y + 3 \geqslant 0,\\
2x + y - 3 \leqslant 0,
\end{cases}\) 则目标函数 \(z = x + 6y\) 的最大值为（\quad）

\choices
    {3}
    {4}
    {18}
    {40}
\end{problem}

\begin{problem}
阅读程序框图，运行相应的程序，则输出 S 的值为（\quad）

\texfigure[max width=6.0cm]{img_repaint/2015/tianjin_science/q3_fig1.tex}

\choices
    {\(-10\)}
    {6}
    {14}
    {18}
\end{problem}

\begin{problem}
设 \(x \in \R\)，则“\(|x - 2| < 1\)”是“\(x^2 + x - 2 > 0\)”的（\quad）

\choices
    {充分而不必要条件}
    {必要而不充分条件}
    {充要条件}
    {既不充分也不必要条件}
\end{problem}

\begin{problem}
如图，在圆 \(O\) 中，\(M\)，\(N\) 是弦 \(AB\) 的三等分点，弦 \(CD\)，\(CE\) 分别经过点 \(M\)，\(N\)，若 \(CM = 2\)，\(MD = 4\)，\(CN = 3\)，则线段 \(NE\) 的长为（\quad）

\bitmapfigure[max width=0.42\linewidth]{img/2015/tianjin_science/q5_fig1.png}

\choices
    {\(\frac{8}{3}\)}
    {3}
    {\(\frac{10}{3}\)}
    {\(\frac{5}{2}\)}
\end{problem}

\begin{problem}
已知双曲线 \(\frac{x^2}{a^2} - \frac{y^2}{b^2} = 1\)（\(a > 0, b > 0\)）的一条渐近线过点 \((2, \sqrt{3})\)，且双曲线的一个焦点在抛物线 \(y^2 = 4\sqrt{7} x\) 的准线上，则双曲线的方程为（\quad）

\choices
    {\(\frac{x^2}{21} -\frac{y^2}{28} = 1\)}
    {\(\frac{x^2}{28} -\frac{y^2}{21} = 1\)}
    {\(\frac{x^2}{9} -\frac{y^2}{4} = 1\)}
    {\(\frac{x^2}{4} -\frac{y^2}{3} = 1\)}
\end{problem}

\begin{problem}
已知定义在 \(\R\) 上的函数 \(f(x) = 2^{|x - m|} - 1\) （\(m\) 为实数） 为偶函数，记 \(a = f(\log_{0.5} 3)\)，\(b = f(\log_{2} 5)\)，\(c = f(2m)\)，则 \(a\)，\(b\)，\(c\) 的大小关系为（\quad）

\choices
    {\(a < b < c\)}
    {\(a < c < b\)}
    {\(c < a < b\)}
    {\(c < b < a\)}
\end{problem}

\begin{problem}
已知函数 \( f(x) = \begin{cases}
2 - |x|, & x \leqslant 2,\\
(x - 2)^2, & x > 2,
\end{cases} \) 函数 \( g(x) = b - f(2 - x) \)，其中 \( b \in \R \). 若函数 \( y = f(x) - g(x) \) 恰有4个零点，则 \( b \) 的取值范围是（\quad）

\choices
    {\(\left(\frac{7}{4}, +\infty\right)\)}
    {\(\left(-\infty, \frac{7}{4}\right)\)}
    {\(\left(0, \frac{7}{4}\right)\)}
    {\(\left(\frac{7}{4}, 2\right)\)}
\end{problem}

\section{填空题}

\begin{problem}
\(\i\) 是虚数单位，若复数 \((1 - 2\i)(a + \i)\) 是纯虚数，则实数 \(a\) 的值为\fillinblank{}.
\end{problem}

\begin{problem}
一个几何体的三视图如图所示（单位：\(\mathrm{m}\)），则该几何体的体积为\fillinblank{}\(\mathrm{m}^3\).
\bitmapfigure[max width=0.58\linewidth]{img/2015/tianjin_science/q10_fig1.png}
\end{problem}

\begin{problem}
曲线 \( y = x^{2} \) 与直线 \( y = x \) 所围成的封闭图形的面积为\fillinblank{}.
\end{problem}

\begin{problem}
在 \(\left(x - \frac{1}{4x}\right)^6\) 的展开式中，\(x^2\) 的系数为\fillinblank{}.
\end{problem}

\begin{problem}
在 \(\triangle ABC\) 中，内角 \(A\)，\(B\)，\(C\) 所对的边分别为 \(a\)，\(b\)，\(c\). 已知 \(\triangle ABC\) 的面积为 \(3\sqrt{15}\)，\(b - c = 2\)，\(\cos A = -\frac{1}{4}\)，则 \(a\) 的值为\fillinblank{}.
\end{problem}

\begin{problem}
在等腰梯形 \(ABCD\) 中，已知 \(AB \parallel DC\)，\(AB = 2\)，\(BC = 1\)，\(\angle ABC = 60^\circ\). 动点 \(E\) 和 \(F\) 分别在线段 \(BC\) 和 \(DC\) 上，且 \(\overrightarrow{BE} = \lambda \overrightarrow{BC}\)，\(\overrightarrow{DF} = \frac{1}{9}\overrightarrow{DC}\)，则 \(\overrightarrow{AE} \cdot \overrightarrow{AF}\) 的最小值为\fillinblank{}.
\end{problem}

\section{解答题}

\begin{problem}
已知函数 \(f(x) = \sin^2 x - \sin^2 \left( x - \frac{\pi}{6} \right)\)，\(x \in \R\).

\begin{enumerate}
\item 求 \( f(x) \) 的最小正周期.
\item 求 \( f(x) \) 在区间 \( \left[-\frac{\pi}{3}, \frac{\pi}{4}\right] \) 上的最大值和最小值.
\end{enumerate}
\end{problem}

\begin{problem}
为推动乒乓球运动的发展，某乒乓球比赛允许不同协会的运动员组队参加. 现有来自甲协会的运动员3名，其中种子选手2名；乙协会的运动员5名，其中种子选手3名. 从这8名运动员中随机选择4人参加比赛.

\begin{enumerate}
\item 设 A 为事件“选出的 4 人中恰有 2 名种子选手，且这 2 名种子选手来自同一个协会”，求事件 A 发生的概率；
\item 设 X 为选出的 4 人中种子选手的人数，求随机变量 X 的分布列和数学期望.
\end{enumerate}
\end{problem}

\begin{problem}
如图，在四棱柱 \(ABCD - A_{1}B_{1}C_{1}D_{1}\) 中，侧棱 \(A_{1}A \perp\) 底面 \(ABCD\)，\(AB \perp AC\)，\(AB = 1\)，\(AC = AA_{1} = 2\)，\(AD = CD = \sqrt{5}\)，且点 \(M\) 和 \(N\) 分别为 \(B_{1}C\) 和 \(D_{1}D\) 的中点.

\begin{enumerate}
\item 求证： \(MN \parallel\) 平面 \(ABCD\)；
\item 求二面角 \(D_{1} - AC - B_{1}\) 的正弦值；
\item 设 \(E\) 为棱 \(A_{1}B_{1}\) 上的点，若直线 \(NE\) 和平面 \(ABCD\) 所成角的正弦值为 \(\frac{1}{3}\)，求线段 \(A_{1}E\) 的长.
\end{enumerate}
\end{problem}

\begin{problem}
已知数列 \(\{a_{n}\}\) 满足 \(a_{n + 2} = qa_{n}\) （\(q\) 为实数，且 \(q \neq 1\)）， \(n \in \mathbf{N}^{*}\)，\(a_{1} = 1\)，\(a_{2} = 2\)，且 \(a_{2} + a_{3}\)，\(a_{3} + a_{4}\)，\(a_{4} + a_{5}\)，成等差数列.

\begin{enumerate}
\item 求 q 的值和 \( \{a_{n}\} \) 的通项公式；
\item 设 \(b_{n}=\frac{\log_{2}a_{2n}}{a_{2n-1}}\)，\(n\in \mathbf{N}^{*}\)，求数列 \( \{b_{n}\} \) 的前 n 项和.
\end{enumerate}
\end{problem}

\begin{problem}
已知椭圆 \(\frac{x^2}{a^2} + \frac{y^2}{b^2} = 1\)（\(a > b > 0\)）的左焦点为 \(F(-c, 0)\)，离心率为 \(\frac{\sqrt{3}}{3}\)，点 \(M\) 在椭圆上且位于第一象限，直线 \(FM\) 被圆 \(x^2 + y^2 = \frac{b^2}{4}\) 截得的线段的长为 \(c\)，\(|FM| = \frac{4\sqrt{3}}{3}\).

\begin{enumerate}
\item 求直线 \(FM\) 的斜率；
\item 求椭圆的方程；
\item 设动点 \(P\) 在椭圆上，若直线 \(FP\) 的斜率大于 \(\sqrt{2}\)，求直线 \(OP\) （\(O\) 为原点） 的斜率的取值范围.
\end{enumerate}
\end{problem}

\begin{problem}
已知函数 \(f(x) = nx - x^n\)，\(x \in \R\). 其中 \( n \in \mathbf{N}^* \)，且 \( n \geqslant 2 \).

\begin{enumerate}
\item 讨论 \( f(x) \) 的单调性；
\item 设曲线 \( y = f(x) \) 与 \( x \) 轴正半轴的交点为 \(P\)，曲线在点 \(P\) 处的切线方程为 \( y = g(x) \). 求证：对于任意的正实数 \( x \)，都有 \( f(x) \leqslant g(x) \)；
\item 若关于 \(x\) 的方程 \(f(x)=a\)（\(a\) 为实数）有两个正实数根 \(x_{1}\)，\(x_{2}\)，求证： \(|x_{2} - x_{1}| < \frac{a}{1 - n} + 2\).
\end{enumerate}
\end{problem}
